% GENERATED by gen_appA_r4.py from appA_template_r4.tex (Lean text from zeta-23-lean-release, release-v1 at 8c15295801687f63bd3a8c5c96822800565303cf). Do not edit; edit the template and re-run.
\section{The Lean statements}\label{app:lean}
\sloppy

The Lean text below is copied verbatim, by a script that reads the files, from the release of the formalisation: the repository
\repourl{} at tag \repotag{} (\S\ref{app:layout}). Only line breaks inside long lines are added by the typesetting.

\emph{What is claimed} is stated in three files, which import nothing but Mathlib and one another: \texttt{comparator/ChallengeDeps.lean} of the Lean artifact
of~\cite{AF}, which defines the zeros of \texttt{riemannZeta} and the counting functions (\S\ref{app:counting}); the definition layer of this
paper, \texttt{comparator/ChallengeDeps/FollowUpZeta.lean}, with the windows, the local inequalities and the three certificate
\texttt{Prop}s (\S\ref{app:defs}); and the fourteen statements of record, \texttt{comparator/Challenge/FollowUpZeta.lean}
(\S\ref{app:record}). A reader who wants to know what is claimed needs only these three files. The statements of record carry
\texttt{sorry} by design: they are the challenge side of the comparator set-up of the repository of~\cite{AF}. The file
\texttt{comparator/Solution/FollowUpZeta.lean} proves exactly these fourteen statements from the library \texttt{ZetaS}, whose definitions
are copies of those of \S\ref{app:defs} in the namespace \texttt{ZetaS}; the conversions between the two are proofs, not trusted text. The
library's theorems are in \S\ref{app:headline}, the chain for the $K=5$ certificate in \S\ref{app:chain}, the replay's theorems in
\S\ref{app:replay} and the majorant in \S\ref{app:majorant}.

The doc-strings are part of the Lean text and are reproduced unchanged. They are development notes, not part of what is claimed, and
not all current: they refer to an earlier draft of this paper by its labels (for instance \texttt{thm:zeta-mainw2}
is Theorem~\ref{thm:S}, \texttt{thm:sigd-Sigma} and \texttt{thm:sigd-D} are Theorem~\ref{thm:main}, \texttt{cor:oll-SC} is
Corollary~\ref{cor:SC}, \texttt{lem:sigd-env} is Lemma~\ref{lem:room}) and to its line numbers, they name work packages, folders and
dates of the development, and ``level A'' in them means kernel-checked in the sense of \S\ref{ssec:formal}.
The data files that the doc-strings name are in \texttt{supplementary/certificates/}, with the same sha256 (Table~\ref{tab:certs}). The
header comment of \S4 of the definition layer says that the certificates were established in Arb ball arithmetic; as \S\ref{ssec:bnb} states,
the box bounds were computed in outward-rounded float64 intervals, with the kernel values enclosed in Arb. Only the Lean code is claimed.

\subsection{The counting functions}\label{app:counting}
From \texttt{comparator/ChallengeDeps.lean} of the repository of~\cite{AF}, unchanged from the tree as shipped (the functions used below are
\texttt{Ncount}, \texttt{Ndist}, \texttt{N0simple}, \texttt{Nsimple} and \texttt{zerosIn}). The phrase ``Not used by the challenge
statements'' in two of these doc-strings refers to the challenge statements of~\cite{AF}; \texttt{Nsimple} is used by those of
\S\ref{app:record}.

\begin{Verbatim}[fontsize=\footnotesize,breaklines,frame=single,framesep=2mm]
/-! ## 1. Nontrivial zeros of ζ and the counting functions (Theorems A–D) -/

/-- ρ is a nontrivial zero of the Riemann zeta function: ζ(ρ) = 0 with 0 < Re ρ < 1 (the open
critical strip). Every zero in the strip is "nontrivial" in the usual sense (it is neither a trivial
zero −2(n+1), of real part ≤ −2, nor the pole s = 1); the converse — that every zero other than the
trivial ones lies in the open strip — is classical and is not needed to state the theorems. -/
def IsNontrivialZero (ρ : ℂ) : Prop := riemannZeta ρ = 0 ∧ 0 < ρ.re ∧ ρ.re < 1

/-- m_ρ, the multiplicity of ρ: the order of vanishing of ζ at ρ, via Mathlib's `analyticOrderAt`
(ℕ∞-valued; `toNat` sends ⊤ to 0, so `1 ≤ zeroMult ρ` encodes both "ζ is not locally identically
zero at ρ" and "ρ is a zero"). -/
def zeroMult (ρ : ℂ) : ℕ := (analyticOrderAt riemannZeta ρ).toNat

/-- The nontrivial zeros with ordinate in the window (T₁, T₂]: {ρ | ζ(ρ) = 0, 0 < Re ρ < 1,
T₁ < Im ρ ≤ T₂}. (Positive-ordinate window, not |Im ρ|.) -/
def zerosIn (T₁ T₂ : ℝ) : Set ℂ := {ρ | IsNontrivialZero ρ ∧ T₁ < ρ.im ∧ ρ.im ≤ T₂}

/-- N(T₁,T₂): the number of nontrivial zeros with T₁ < Im ρ ≤ T₂, counted WITH multiplicity
(`∑ᶠ` is Mathlib's finite sum; that the window contains finitely many zeros is proved on the
solution side, not assumed). -/
def Ncount (T₁ T₂ : ℝ) : ℕ := ∑ᶠ ρ ∈ zerosIn T₁ T₂, zeroMult ρ

/-- N_d(T₁,T₂): the number of nontrivial zeros with T₁ < Im ρ ≤ T₂, each distinct zero counted
once. -/
def Ndist (T₁ T₂ : ℝ) : ℕ := (zerosIn T₁ T₂).ncard

/-- N₀(T₁,T₂): the zeros ON the critical line Re ρ = 1/2 with T₁ < Im ρ ≤ T₂, with multiplicity.
(Not used by the challenge statements; kept so that the block matches Zeta23/Statement.lean.) -/
def N0 (T₁ T₂ : ℝ) : ℕ := ∑ᶠ ρ ∈ zerosIn T₁ T₂ ∩ {ρ | ρ.re = 1 / 2}, zeroMult ρ

/-- N₀*(T₁,T₂): the number of DISTINCT zeros ON the critical line Re ρ = 1/2 with
T₁ < Im ρ ≤ T₂. This is what Theorems A and D bound from below. -/
def N0star (T₁ T₂ : ℝ) : ℕ := (zerosIn T₁ T₂ ∩ {ρ | ρ.re = 1 / 2}).ncard

/-- N₀ˢ(T₁,T₂): the number of SIMPLE zeros (multiplicity exactly 1) ON the critical line with
T₁ < Im ρ ≤ T₂. -/
def N0simple (T₁ T₂ : ℝ) : ℕ := (zerosIn T₁ T₂ ∩ {ρ | ρ.re = 1 / 2} ∩ {ρ | zeroMult ρ = 1}).ncard

/-- Nˢ(T₁,T₂): the number of simple zeros with T₁ < Im ρ ≤ T₂ (anywhere in the strip).
(Not used by the challenge statements; kept so that the block matches Zeta23/Statement.lean.) -/
def Nsimple (T₁ T₂ : ℝ) : ℕ := (zerosIn T₁ T₂ ∩ {ρ | zeroMult ρ = 1}).ncard
\end{Verbatim}

\subsection{The definitions of this paper}\label{app:defs}
The definition layer \texttt{comparator/ChallengeDeps/FollowUpZeta.lean} (namespace \texttt{FollowUpZeta}), from its first line of code to
its end; the file's header comment, which lists its contents, is omitted.

\begin{Verbatim}[fontsize=\footnotesize,breaklines,frame=single,framesep=2mm]
import ChallengeDeps

open scoped BigOperators
open Finset

noncomputable section

namespace FollowUpZeta

/-! ## §1 Zeros that are simple or on the critical line -/

/-- N^sc(T₁,T₂) (cor:oll-SC): the nontrivial zeros with T₁ < Im ρ ≤ T₂ that lie ON the critical line, counted WITH
multiplicity, plus those OFF the line that are simple (both members ρ, 1 − ρ̄ of each simple off-line pair; multiple
off-line zeros are not counted). An off-line zero in the set is simple, so its weight `zeroMult ρ` is 1. -/
def Nsc (T₁ T₂ : ℝ) : ℕ :=
  ∑ᶠ ρ ∈ zerosIn T₁ T₂ ∩ ({ρ | ρ.re = 1 / 2} ∪ {ρ | zeroMult ρ = 1}), zeroMult ρ

/-! ## §2 Windows and the normalised kernel -/

/-- The polynomial window ψ̃ of thm:zeta-mainw2:
ψ̃(s) = (50000/50037)(1 + (127/2500)u − (8997/10000)u² + (6187/10000)u³ + (3/1250)u⁴), u = (2s)²,
on [−1/2, 1/2] (values outside are never read). -/
def psiPoly8A (s : ℝ) : ℝ :=
  50000 / 50037 * (1 + 127 / 2500 * (2 * s) ^ 2 - 8997 / 10000 * ((2 * s) ^ 2) ^ 2
    + 6187 / 10000 * ((2 * s) ^ 2) ^ 3 + 3 / 1250 * ((2 * s) ^ 2) ^ 4)

/-- The window ψ(s) = cos(1.6 s) of thm:sigd-Sigma and thm:sigd-D. -/
def psiCos16 (s : ℝ) : ℝ := Real.cos (8 / 5 * s)

/-- The normalised Fourier transform of a window: k_ψ(x) = ∫_{−1/2}^{1/2} ψ(s) cos(2πxs) ds / ∫_{−1/2}^{1/2} ψ. -/
def kPsi (ψ : ℝ → ℝ) (x : ℝ) : ℝ :=
  (∫ s in (-(1 / 2 : ℝ))..(1 / 2), ψ s * Real.cos (2 * Real.pi * x * s))
    / ∫ s in (-(1 / 2 : ℝ))..(1 / 2), ψ s

/-! ## §3 The local inequalities -/

/-- Position weights for K consecutive points (lem:zeta-Dprime), 0-based: pair weights `γ s i` for spans
1 ≤ s ≤ K − 1 and positions 0 ≤ i < K − s, nonnegative and summing to 2 for every span; gap penalties `μ l ≥ 0`,
l : Fin (K − 1). Values of `γ` outside the index range are never read. -/
structure LocalWeights (K : ℕ) where
  γ : ℕ → ℕ → ℝ
  μ : Fin (K - 1) → ℝ
  two_le : 2 ≤ K
  γ_nonneg : ∀ s ∈ Icc 1 (K - 1), ∀ i ∈ range (K - s), 0 ≤ γ s i
  γ_sum : ∀ s ∈ Icc 1 (K - 1), ∑ i ∈ range (K - s), γ s i = 2
  μ_nonneg : ∀ l, 0 ≤ μ l

/-- ν := Σ_i μ_i. -/
def LocalWeights.nu {K : ℕ} (W : LocalWeights K) : ℝ := ∑ l, W.μ l

/-- The span g_i + ⋯ + g_{i+s−1} of a gap vector (0-based; indices ≥ K − 1 contribute nothing). -/
def gapSpan {K : ℕ} (g : Fin (K - 1) → ℝ) (i s : ℕ) : ℝ :=
  ∑ l ∈ univ.filter (fun l : Fin (K - 1) => i ≤ l.val ∧ l.val < i + s), g l

/-- F_{γ,μ}(g) = Σ_i μ_i g_i + Σ_{s=1}^{K−1} Σ_i γ_{s,i} k(g_i + ⋯ + g_{i+s−1})² (eq:zeta-LIw). -/
def localF (k : ℝ → ℝ) {K : ℕ} (W : LocalWeights K) (g : Fin (K - 1) → ℝ) : ℝ :=
  (∑ l, W.μ l * g l)
    + ∑ s ∈ Icc 1 (K - 1), ∑ i ∈ range (K - s), W.γ s i * k (gapSpan g i s) ^ 2

/-- (LI′) of lem:zeta-Dprime: F_{γ,μ}(g) ≥ c for every gap vector g ∈ [0,∞)^{K−1}. -/
def LocalCert (k : ℝ → ℝ) {K : ℕ} (W : LocalWeights K) (c : ℝ) : Prop :=
  ∀ g : Fin (K - 1) → ℝ, (∀ l, 0 ≤ g l) → c ≤ localF k W g

/-- All-marks data (def:zeta-LIm): position weights as in `LocalWeights` and site credits `b i j` for site i : Fin K and
capped mark j : Fin 2 (j = 0 ↦ mark 1, j = 1 ↦ mark 2); no sign condition. -/
structure MarkWeights (K : ℕ) extends LocalWeights K where
  b : Fin K → Fin 2 → ℝ

/-- a_j := Σ_i b_i(j); `a 0 = a₁`, `a 1 = a₂`. -/
def MarkWeights.a {K : ℕ} (W : MarkWeights K) (j : Fin 2) : ℝ := ∑ i, W.b i j

/-- The numerical mark m_i ∈ {1, 2} at 0-based position i of a pattern (0 outside [0, K)). -/
def markVal {K : ℕ} (m : Fin K → Fin 2) (i : ℕ) : ℝ :=
  if h : i < K then ((m ⟨i, h⟩ : ℕ) : ℝ) + 1 else 0

/-- F_𝐦(g) = Σ_{s,i} γ_{s,i} m_i m_{i+s} k(g_i + ⋯ + g_{i+s−1})² + Σ_i μ_i g_i (def:zeta-LIm). -/
def localFm (k : ℝ → ℝ) {K : ℕ} (W : MarkWeights K) (m : Fin K → Fin 2) (g : Fin (K - 1) → ℝ) : ℝ :=
  (∑ s ∈ Icc 1 (K - 1), ∑ i ∈ range (K - s),
      W.γ s i * markVal m i * markVal m (i + s) * k (gapSpan g i s) ^ 2)
    + ∑ l, W.μ l * g l

/-- (LI_m) of def:zeta-LIm: for EVERY mark pattern 𝐦 ∈ {1,2}^K and every g ∈ [0,∞)^{K−1}, F_𝐦(g) ≥ Σ_i b_i(m_i). -/
def LocalCertAM (k : ℝ → ℝ) {K : ℕ} (W : MarkWeights K) : Prop :=
  ∀ m : Fin K → Fin 2, ∀ g : Fin (K - 1) → ℝ, (∀ l, 0 ≤ g l) → ∑ i, W.b i (m i) ≤ localFm k W m g

/-! ## §4 The certificate hypotheses

Each Prop asserts that SOME weights with the stated constants satisfy the local inequality on the whole orthant. The
weights that realise it are exact rationals in the data files named below (identified by sha256); the inequality for those
weights was established outside Lean by interval branch-and-bound (Arb ball arithmetic, every box accepted only when a
rigorous lower bound exceeds the claim; the unbounded part of the orthant is discarded exactly, since
F ≥ (min μ)·Σ g there). Paths are relative to the development's project folder; the files named here are in
`supplementary/certificates/` of the repository, some logs under new names (map: `supplementary/README.md`). -/

/-- **Certificate S8** (thm:zeta-mainw2, rem:zeta-P8, lem:zeta-Dprime): K = 8, window ψ̃, weights with ν = 7/1700 for
which (LI′) holds at the claim c = 0.00796.
Data: `round2/d6_5_numerics/w_poly8A_K8_mu1700.json`, sha256
`257b340d9da4f343600a21589f16d8915f3e082499d475e153c5453980671e78`; window `round2/d6_5_numerics/win_poly8A.json`,
sha256 `c4005f57a7920f78aa2df33b093c9c5030efac2b8e39fc6cf7371b0c0f32dc14` (the unnormalised v = (50037/50000)ψ̃; the kernel
is homogeneous of degree 0). Certifier `round2/d6_5_numerics/bnb_lb3.py` with `polywin.py`; runs
`round2/d6_5_numerics/runs/P8/DRIVER.log` and the independent re-run `runs/P8b/DRIVER.log`, both ending
`ALL SHARDS OK: certificate claim=0.00796 holds` (about 4.5·10⁸ boxes each). -/
def CertS8 : Prop :=
  ∃ W : LocalWeights 8, W.nu = 7 / 1700 ∧ LocalCert (kPsi psiPoly8A) W (199 / 25000)

/-- **Certificate AM5** (thm:zeta-allmarks, rem:sigd-cert, the K = 5 row): all-marks data with a₁ = 1280197/10⁸,
a₂ = 48749/3125000, ν = 1/125 for which (LI_m) holds for k_ψ, ψ = cos 1.6s.
Data: `round1/X2_numerics/claims_a1.6_K5_mu500.json` (byte-identical to `claims_d15e6_a1.6_K5_mu500.json`), sha256
`cf055dd2ee9a9e469b6510fe70b5fbc416b8a40515cc6746bc390d40d88777f2`; the 20 class files `xpat_K5_*.json` of the same
folder, sha256 of their concatenation in sorted file-name order
`a5947c01d34ddcbc7366861b35e83942099fcbf7ac71e13fcd432796882dd34e`. Certifier `round1/X2_numerics/bnb_interval_w.py`;
logs `cert_K5_w_a16_mu500_d15e6.txt` and `cert_K5_w_a16_mu500_d15e6_b.txt` (20/20 reversal classes ok, covering the 32
patterns; the weights are reversal-symmetric). A replay of this certificate by the Lean kernel was completed on
28 September 2026: `ZetaS.CertV2.certAM5_replayed` proves the library's copy `ZetaS.CertAM5`, and the hypothesis-free
K = 5 theorems are in the module `ReplayHeadlines` of the library `ZetaSReplay`. In the statements of this topic the
Prop remains a displayed hypothesis, like the other two. -/
def CertAM5 : Prop :=
  ∃ W : MarkWeights 5, W.a 0 = 1280197 / 10 ^ 8 ∧ W.a 1 = 48749 / 3125000 ∧ W.nu = 1 / 125 ∧
    LocalCertAM (kPsi psiCos16) W

/-- **Certificate AM7** (thm:sigd-Sigma, thm:sigd-D, rem:sigd-cert, the K = 7 row): all-marks data with
a₁ = 1824837/10⁸, a₂ = 1168069/(5·10⁷), ν = 3/250 for which (LI_m) holds for k_ψ, ψ = cos 1.6s.
Data: `round2/d8_9_numerics/k7m500/rerun/claims_a1.6_K7_mu500_L1e5.json` (the final claims, every claim of the LP lowered by
exactly 10⁻⁵), sha256 `d2cf393019407bdc797c73147ec8e3dc66ec41fe745600c84173f74ac66bfc5c`; the 72 class files
`xpat_K7_*.json` of the same folder, sha256 of their concatenation in sorted file-name order
`33c4c8c5762bcebdf8f544ab34a137a6a7ba4065d08f5caee0a51188e1a67514`. Certifier `round2/d8_9_numerics/bnb_lb3_w.py`
(driver `runlb3.sh`); certificate log `round2/d8_9_numerics/k7m500/rerun/cert_lb3_rerun.txt` (72/72 reversal classes ok,
each at exactly the final claim); completeness (the 72 classes cover all 128 patterns, each certified at a claim ≥ the
required one) `round2/d8_21_numerics/k7_verify_pexact2.txt`. The weights are reversal-symmetric. -/
def CertAM7 : Prop :=
  ∃ W : MarkWeights 7, W.a 0 = 1824837 / 10 ^ 8 ∧ W.a 1 = 1168069 / (5 * 10 ^ 7) ∧ W.nu = 3 / 250 ∧
    LocalCertAM (kPsi psiCos16) W

end FollowUpZeta
\end{Verbatim}

\subsection{The fourteen statements of record}\label{app:record}
The statement file \texttt{comparator/Challenge/FollowUpZeta.lean}, from its first line of code to its end (the header comment is omitted).
The six headline statements are those of Theorems~\ref{thm:main} and~\ref{thm:S} and Corollary~\ref{cor:SC}, each with its single
certificate hypothesis; the six \texttt{\_cumulative} forms are the same bounds on the windows $0<\Im\rho\le T$; the last two statements
show that the counts are genuine (every window holds finitely many zeros) and that the denominator tends to infinity.

\begin{Verbatim}[fontsize=\footnotesize,breaklines,frame=single,framesep=2mm]
import ChallengeDeps.FollowUpZeta

open FollowUpZeta

noncomputable section

/-- **Simple zeros (K = 5 certificate)**, thm:sigd-Sigma with the K = 5 row of rem:sigd-cert: at least 0.675158622 of
the nontrivial zeros with T < Im ρ ≤ 2T (counted with multiplicity) are simple, for all large T. -/
theorem zeta_simple_K5 (hcert : CertAM5) :
    ∀ ε > 0, ∃ T₀ : ℝ, ∀ T ≥ T₀,
      ((675158622 : ℝ) / 10 ^ 9 - ε) * (Ncount T (2 * T) : ℝ) ≤ Nsimple T (2 * T) := by
  sorry

/-- `zeta_simple_K5` on the cumulative windows 0 < Im ρ ≤ T (same constant). -/
theorem zeta_simple_K5_cumulative (hcert : CertAM5) :
    ∀ ε > 0, ∃ T₀ : ℝ, ∀ T ≥ T₀,
      ((675158622 : ℝ) / 10 ^ 9 - ε) * (Ncount 0 T : ℝ) ≤ Nsimple 0 T := by
  sorry

/-- **Distinct zeros (K = 5 certificate)**, thm:sigd-D with the K = 5 row of rem:sigd-cert: at least 0.837579311. -/
theorem zeta_distinct_K5 (hcert : CertAM5) :
    ∀ ε > 0, ∃ T₀ : ℝ, ∀ T ≥ T₀,
      ((837579311 : ℝ) / 10 ^ 9 - ε) * (Ncount T (2 * T) : ℝ) ≤ Ndist T (2 * T) := by
  sorry

/-- `zeta_distinct_K5` on the cumulative windows 0 < Im ρ ≤ T (same constant). -/
theorem zeta_distinct_K5_cumulative (hcert : CertAM5) :
    ∀ ε > 0, ∃ T₀ : ℝ, ∀ T ≥ T₀,
      ((837579311 : ℝ) / 10 ^ 9 - ε) * (Ncount 0 T : ℝ) ≤ Ndist 0 T := by
  sorry

/-- **Simple zeros**, thm:sigd-Sigma (K = 7 certificate): at least 0.676102666. -/
theorem zeta_simple_K7 (hcert : CertAM7) :
    ∀ ε > 0, ∃ T₀ : ℝ, ∀ T ≥ T₀,
      ((676102666 : ℝ) / 10 ^ 9 - ε) * (Ncount T (2 * T) : ℝ) ≤ Nsimple T (2 * T) := by
  sorry

/-- `zeta_simple_K7` on the cumulative windows 0 < Im ρ ≤ T (same constant). -/
theorem zeta_simple_K7_cumulative (hcert : CertAM7) :
    ∀ ε > 0, ∃ T₀ : ℝ, ∀ T ≥ T₀,
      ((676102666 : ℝ) / 10 ^ 9 - ε) * (Ncount 0 T : ℝ) ≤ Nsimple 0 T := by
  sorry

/-- **Distinct zeros**, thm:sigd-D (K = 7 certificate): at least 0.838051333. -/
theorem zeta_distinct_K7 (hcert : CertAM7) :
    ∀ ε > 0, ∃ T₀ : ℝ, ∀ T ≥ T₀,
      ((838051333 : ℝ) / 10 ^ 9 - ε) * (Ncount T (2 * T) : ℝ) ≤ Ndist T (2 * T) := by
  sorry

/-- `zeta_distinct_K7` on the cumulative windows 0 < Im ρ ≤ T (same constant). -/
theorem zeta_distinct_K7_cumulative (hcert : CertAM7) :
    ∀ ε > 0, ∃ T₀ : ℝ, ∀ T ≥ T₀,
      ((838051333 : ℝ) / 10 ^ 9 - ε) * (Ncount 0 T : ℝ) ≤ Ndist 0 T := by
  sorry

/-- **Simple zeros on the critical line**, thm:zeta-mainw2 (window ψ̃, K = 8): at least 0.6733736895. -/
theorem zeta_simple_on_line (hcert : CertS8) :
    ∀ ε > 0, ∃ T₀ : ℝ, ∀ T ≥ T₀,
      ((6733736895 : ℝ) / 10 ^ 10 - ε) * (Ncount T (2 * T) : ℝ) ≤ N0simple T (2 * T) := by
  sorry

/-- `zeta_simple_on_line` on the cumulative windows 0 < Im ρ ≤ T (same constant). -/
theorem zeta_simple_on_line_cumulative (hcert : CertS8) :
    ∀ ε > 0, ∃ T₀ : ℝ, ∀ T ≥ T₀,
      ((6733736895 : ℝ) / 10 ^ 10 - ε) * (Ncount 0 T : ℝ) ≤ N0simple 0 T := by
  sorry

/-- **Zeros that are simple or on the critical line**, cor:oll-SC (on-line zeros with multiplicity plus simple
off-line zeros): at least 0.8879195. -/
theorem zeta_simple_or_critical (hcert : CertS8) :
    ∀ ε > 0, ∃ T₀ : ℝ, ∀ T ≥ T₀,
      ((8879195 : ℝ) / 10 ^ 7 - ε) * (Ncount T (2 * T) : ℝ) ≤ Nsc T (2 * T) := by
  sorry

/-- `zeta_simple_or_critical` on the cumulative windows 0 < Im ρ ≤ T (same constant). -/
theorem zeta_simple_or_critical_cumulative (hcert : CertS8) :
    ∀ ε > 0, ∃ T₀ : ℝ, ∀ T ≥ T₀,
      ((8879195 : ℝ) / 10 ^ 7 - ε) * (Ncount 0 T : ℝ) ≤ Nsc 0 T := by
  sorry

/-- Non-vacuity companion: every window T₁ < Im ρ ≤ T₂ holds finitely many nontrivial zeros, so `Ncount`, `Nsc`
(finite sums) and `Nsimple`, `Ndist`, `N0simple` (`Set.ncard`) are genuine counts (Lean's `finsum` and `ncard` would
return 0 on an infinite set). -/
theorem zeta_zerosIn_finite (T₁ T₂ : ℝ) :
    (zerosIn T₁ T₂).Finite := by
  sorry

/-- Non-vacuity companion: the denominator N(T, 2T) tends to infinity. -/
theorem zeta_tendsto_Ncount :
    Filter.Tendsto (fun T : ℝ => (Ncount T (2 * T) : ℝ)) Filter.atTop Filter.atTop := by
  sorry

end
\end{Verbatim}

\subsection{The library's headline theorems}\label{app:headline}
The four certificate-only theorems for simple and distinct zeros (\texttt{ZetaS/Top/TopFinal.lean}). They apply the forms with the majorant
threshold $\frac{33}{10}$ to the proved majorant certificate \texttt{MajSound.certMajV2\_proved} (\S\ref{app:majorant}), so that their only
hypothesis is the local certificate:

\begin{Verbatim}[fontsize=\footnotesize,breaklines,frame=single,framesep=2mm]
namespace ZetaS
namespace Top

/-- `zeta_simple_K5_V2` with the majorant certificate proved (`MajSound.certMajV2_proved`). -/
theorem zeta_simple_K5_final (hcert : CertAM5) :
    ∀ ε > 0, ∃ T₀ : ℝ, ∀ T ≥ T₀,
      ((675158622 : ℝ) / 10 ^ 9 - ε) * (Ncount T (2 * T) : ℝ) ≤ Nsimple T (2 * T) :=
  zeta_simple_K5_V2 hcert MajSound.certMajV2_proved

/-- `zeta_distinct_K5_V2` with the majorant certificate proved (`MajSound.certMajV2_proved`). -/
theorem zeta_distinct_K5_final (hcert : CertAM5) :
    ∀ ε > 0, ∃ T₀ : ℝ, ∀ T ≥ T₀,
      ((837579311 : ℝ) / 10 ^ 9 - ε) * (Ncount T (2 * T) : ℝ) ≤ Ndist T (2 * T) :=
  zeta_distinct_K5_V2 hcert MajSound.certMajV2_proved

/-- `zeta_simple_K7_V2` with the majorant certificate proved (`MajSound.certMajV2_proved`). -/
theorem zeta_simple_K7_final (hcert : CertAM7) :
    ∀ ε > 0, ∃ T₀ : ℝ, ∀ T ≥ T₀,
      ((676102666 : ℝ) / 10 ^ 9 - ε) * (Ncount T (2 * T) : ℝ) ≤ Nsimple T (2 * T) :=
  zeta_simple_K7_V2 hcert MajSound.certMajV2_proved

/-- `zeta_distinct_K7_V2` with the majorant certificate proved (`MajSound.certMajV2_proved`). -/
theorem zeta_distinct_K7_final (hcert : CertAM7) :
    ∀ ε > 0, ∃ T₀ : ℝ, ∀ T ≥ T₀,
      ((838051333 : ℝ) / 10 ^ 9 - ε) * (Ncount T (2 * T) : ℝ) ≤ Ndist T (2 * T) :=
  zeta_distinct_K7_V2 hcert MajSound.certMajV2_proved

end Top
end ZetaS
\end{Verbatim}

The two theorems for simple zeros on the line and for zeros that are simple or on the line (\texttt{ZetaS/Top/TopSSC.lean}, namespace \texttt{ZetaS.Top});
statements only, each followed in the file by a tactic proof of about ten lines:

\begin{Verbatim}[fontsize=\footnotesize,breaklines,frame=single,framesep=2mm]
/-- **S — simple zeros on the critical line** (the architect's `zeta_simple_on_line`, thm:zeta-mainw2). -/
theorem zeta_simple_on_line (hcert : CertS8) :
    ∀ ε > 0, ∃ T₀ : ℝ, ∀ T ≥ T₀,
      ((6733736895 : ℝ) / 10 ^ 10 - ε) * (Ncount T (2 * T) : ℝ) ≤ N0simple T (2 * T) := by
\end{Verbatim}

\begin{Verbatim}[fontsize=\footnotesize,breaklines,frame=single,framesep=2mm]
/-- **SC — simple or critical zeros** (the architect's `zeta_simple_or_critical`, cor:oll-SC). -/
theorem zeta_simple_or_critical (hcert : CertS8) :
    ∀ ε > 0, ∃ T₀ : ℝ, ∀ T ≥ T₀,
      ((8879195 : ℝ) / 10 ^ 7 - ε) * (Ncount T (2 * T) : ℝ) ≤ Nsc T (2 * T) := by
\end{Verbatim}

The statements of these six theorems are, character for character, those of the six dyadic statements of record in \S\ref{app:record},
with the definitions of the namespace \texttt{ZetaS} in place of those of \texttt{FollowUpZeta}; \texttt{comparator/Solution/FollowUpZeta.lean}
derives the fourteen statements of record from them. The library's definition of \texttt{CertAM5} (\texttt{ZetaS/Top/TopDefs.lean}) is:

\begin{Verbatim}[fontsize=\footnotesize,breaklines,frame=single,framesep=2mm]
/-- **Certificate AM5**: all-marks data with `a₁ = 1280197/10⁸`, `a₂ = 48749/3125000`, `ν = 1/125` for which (LI_m)
holds for `k_ψ`, `ψ = cos 1.6s` (X2's `round1/X2_numerics/xpat_K5_*.json`, `claims_d15e6_a1.6_K5_mu500.json`). -/
def CertAM5 : Prop :=
  ∃ W : MarkWeights 5, W.a 0 = 1280197 / 10 ^ 8 ∧ W.a 1 = 48749 / 3125000 ∧ W.nu = 1 / 125 ∧
    LocalCertAM (kPsi psiCos16) W
\end{Verbatim}

\subsection{The chain for the \texorpdfstring{$K=5$}{K=5} certificate}\label{app:chain}
From \texttt{ZetaS/Cert/ChainV3.lean}: the soundness theorem of the corrected checker and the theorem that twenty checked certificates, one
per canonical class, give \texttt{CertAM5} (the file ends with a diagnostic \texttt{\#eval}, omitted):

\begin{Verbatim}[fontsize=\footnotesize,breaklines,frame=single,framesep=2mm]
namespace ZetaS.CertV2.ChainV3

open ZetaS.CertV2

/-- **The certificate theorem** (node C21, L5_1's proof), against the V3 checker, all dependencies proved. -/
theorem cert_check_sound (C : Cert) (h : C.check = true) : C.data.Holds := Cert.check_sound' C h

/-- **Twenty checked certificates, one per canonical class, give `CertAM5`** (V3 checker, level A). -/
theorem certAM5_of_checks (C : List Cert) (hdata : C.map Cert.data = classes5.map classData)
    (hchk : ∀ c ∈ C, c.check = true) : CertAM5 :=
  certAM5_of_checks_of_sound cert_check_sound C hdata hchk

end ZetaS.CertV2.ChainV3
\end{Verbatim}

\subsection{The replay's theorems: the \texorpdfstring{$K=5$}{K=5} pair without hypothesis}\label{app:replay}
These are in the library \texttt{ZetaSReplay} of the repository, which holds the $1{,}478$ files of the replay (\S\ref{ssec:replay}) in one
flat folder, byte-identical to the files that the kernel checked, and is not among the default build targets; \texttt{lake build
ZetaSReplay} repeats the replay (\S\ref{ssec:reproduce}). The replay compiled all $1{,}478$ files with no failure, in $2$~h~$50$~min of
wall time and $20.4$ Lean-hours, against the library at commit
\texttt{bb7223b}, whose Lean files the release changes only in comments (\S\ref{app:layout}). From \texttt{ZetaSReplay/ReplayAM5.lean}, the
statement of \texttt{certAM5\_replayed} and the first line of its proof (the proof term continues with the twenty theorems
\texttt{R<class>.cert\_ok} proved by the kernel; omitted):

\begin{Verbatim}[fontsize=\footnotesize,breaklines,frame=single,framesep=2mm]
namespace ZetaS.CertV2

theorem certAM5_replayed : ZetaS.CertAM5 :=
  ChainV3.certAM5_of_checks [R11111.cert, R11112.cert, R11121.cert, R11122.cert, R11211.cert, R11212.cert, R11221.cert, R11222.cert, R12112.cert, R12121.cert, R12122.cert, R12212.cert, R12221.cert, R12222.cert, R21112.cert, R21122.cert, R21212.cert, R21222.cert, R22122.cert, R22222.cert] rfl
\end{Verbatim}

From \texttt{ZetaSReplay/ReplayHeadlines.lean}, the two theorems of the $K=5$ pair without hypothesis; their statements are those of
\texttt{ZetaS.Top.zeta\_simple\_K5\_final} and \texttt{zeta\_distinct\_K5\_final} (\S\ref{app:headline}) without the hypothesis
\texttt{hcert}:

\begin{Verbatim}[fontsize=\footnotesize,breaklines,frame=single,framesep=2mm]
namespace ZetaS
namespace Top

/-- **Simple zeros, K = 5, hypothesis-free** (thm:sigd-Sigma with the K = 5 certificate). -/
theorem zeta_simple_K5_uncond :
    ∀ ε > 0, ∃ T₀ : ℝ, ∀ T ≥ T₀,
      ((675158622 : ℝ) / 10 ^ 9 - ε) * (Ncount T (2 * T) : ℝ) ≤ Nsimple T (2 * T) :=
  zeta_simple_K5_final ZetaS.CertV2.certAM5_replayed

/-- **Distinct zeros, K = 5, hypothesis-free** (thm:sigd-D with the K = 5 certificate). -/
theorem zeta_distinct_K5_uncond :
    ∀ ε > 0, ∃ T₀ : ℝ, ∀ T ≥ T₀,
      ((837579311 : ℝ) / 10 ^ 9 - ε) * (Ncount T (2 * T) : ℝ) ≤ Ndist T (2 * T) :=
  zeta_distinct_K5_final ZetaS.CertV2.certAM5_replayed

end Top
end ZetaS
\end{Verbatim}

\subsection{The majorant}\label{app:majorant}
The majorant is not a hypothesis of any headline theorem: it is proved. Its \texttt{Prop} (\texttt{ZetaS/ChallengeZetaS.lean}, namespace
\texttt{ZetaS}), the instance used (\texttt{ZetaS/Majorant/CertMajV2.lean}) and the theorem (\texttt{ZetaS/Majorant/CertMajV2Proof.lean},
namespace \texttt{ZetaS.MajSound}):

\begin{Verbatim}[fontsize=\footnotesize,breaklines,frame=single,framesep=2mm]
/-- **`MajorantCert ψ λ`** — the certified input of lem:sigd-env (l.1283–1296), in the form the proof uses:
an even integrable `B ≥ 0` on `ℝ` with `B ≥ v_ψ := ψ/∫ψ` on `[−1/2,1/2]`, whose Fourier transform
`B̂(t) = ∫ B(s) cos(2πts) ds` vanishes for `|t| ≥ 1`, and with `2 B̂(0) + 4 sup_{[3/4,1]} |B̂| < λ`.
The draft's instance: `B̂` piecewise linear with nodes `j/60`, `β₀ = 1.5556524027186915`, `β* = |β₅₂|`,
and `λ = λ_c = 2 + √(2 − 2a₁)`. -/
def MajorantCert (ψ : ℝ → ℝ) (lam : ℝ) : Prop :=
  ∃ B : ℝ → ℝ, MeasureTheory.Integrable B ∧ (∀ s, B (-s) = B s) ∧ (∀ s, 0 ≤ B s) ∧
    (∀ s ∈ Set.Icc (-(1 / 2 : ℝ)) (1 / 2), ψ s / (∫ t in (-(1 / 2 : ℝ))..(1 / 2), ψ t) ≤ B s) ∧
    (∀ t : ℝ, 1 ≤ |t| → ∫ s, B s * Real.cos (2 * Real.pi * t * s) = 0) ∧
    ∃ βstar : ℝ, (∀ t ∈ Set.Icc (3 / 4 : ℝ) 1, |∫ s, B s * Real.cos (2 * Real.pi * t * s)| ≤ βstar) ∧
      2 * (∫ s, B s) + 4 * βstar < lam
\end{Verbatim}

\begin{Verbatim}[fontsize=\footnotesize,breaklines,frame=single,framesep=2mm]
/-- **CertMajV2** — `MajorantCert psiCos16 (33/10)`. -/
def CertMajV2 : Prop := MajorantCert psiCos16 (33 / 10)
\end{Verbatim}

\begin{Verbatim}[fontsize=\footnotesize,breaklines,frame=single,framesep=2mm]
/-- **CertMajV2**, instantiated with the C05/C06 node statements. -/
theorem certMajV2_proved : CertMajV2 := certMajV2_of_CNodes cnodes_of_statements
\end{Verbatim}

\texttt{ZetaS/Majorant/CertMajProof.lean} also proves, as \texttt{MajSound.certMaj\_proved}, the first version \texttt{CertMaj} of
\texttt{ZetaS/ChallengeZetaS.lean}, with the threshold $3.2063638853$; no theorem uses it.

\subsection{Axiom census and \texttt{sorry}}\label{app:census}
At commit \texttt{bb7223b}, \texttt{\#print axioms} gives exactly \texttt{[propext, Classical.choice, Quot.sound]} for
each of the six theorems of \S\ref{app:headline}, for \texttt{ZetaS.CertV2.ChainV3.certAM5\_of\_checks} and
\texttt{ZetaS.CertV2.ChainV3.cert\_check\_sound}, and for all fourteen theorems of
\texttt{comparator/Solution/FollowUpZeta.lean}. For the replay's three theorems of \S\ref{app:replay} it gives the same three axioms, and
no open declaration lies beneath them. In the repository,
\texttt{comparator/PrintAxioms/FollowUpZeta.lean} prints the axioms of the fourteen theorems of record and
\texttt{comparator/PrintAxioms/FollowUpReplay.lean} those of the three theorems of \S\ref{app:replay} (\S\ref{ssec:reproduce}); the logs and
the \texttt{\#print axioms} outputs of the builds of the release are in \texttt{audit/followup/}.

The library \texttt{ZetaS} has $206$ files and $22{,}751$ lines. It contains $3$ \texttt{sorry}, all in the open node
\texttt{ZetaS/Skeleton/A8g\_SigmaDistGeneral.lean}: Theorems~\ref{thm:OLL} and~\ref{thm:SigmaD} for arbitrary certificate data. That node
is used only by the forms with a general majorant threshold, \texttt{ZetaS.Top.zeta\_simple\_K5\_of\_majorant} and its three companions
(\texttt{ZetaS/Top/TopSolutionGeneral.lean}), and by no headline theorem; the import closure of the fourteen theorems of record contains
no \texttt{sorry}. There is no \texttt{axiom}, \texttt{admit} or \texttt{native\_decide} in the library.

\subsection{Repository layout}\label{app:layout}
The formalisation is published in a modified copy of the Lean repository of~\cite{AF} (\texttt{zeta-23-lean}): \repourl, at tag
\repotag. The commit of the tag is not printed in this paper; it is recorded in the file \texttt{RELEASE\_NOTES.md} of the repository and in
the Zenodo record of the release. The repository is a snapshot of the release without its development history: the commits named in this paper
(such as \texttt{bb7223b}) are commits of that history. The parts that concern this paper:
\begin{itemize}
\item \texttt{ZetaS/} with its root module \texttt{ZetaS.lean}: the library of this paper, outside the default build targets (its folders are
listed below); \texttt{ZetaSCertShims/}: three forwarding modules through which the checker files are imported by their bare names, so
that they are kept byte for byte as checked;
\item \texttt{ZetaSReplay/}: the $1{,}478$ files of the kernel replay of \texttt{CertAM5} ($1{,}475$ generated shards and the three common
modules \texttt{ReplayCommon}, \texttt{ReplayAM5} and \texttt{ReplayHeadlines}), in one flat folder, as a library outside the default build
targets (\S\ref{ssec:replay}, \S\ref{app:replay});
\item \texttt{comparator/}: the comparator set-up of~\cite{AF}, with seven new files for this paper: \texttt{ChallengeDeps/FollowUpZeta.lean} and
\texttt{Challenge/FollowUpZeta.lean} (quoted above), \texttt{Solution/FollowUpZeta.lean}, \texttt{PrintAxioms/FollowUpZeta.lean} and
\texttt{PrintAxioms/FollowUpReplay.lean}, \texttt{config-followup.json} and \texttt{README\_followup.md};
\item \texttt{supplementary/}: the data, certifiers and logs of Table~\ref{tab:certs} and the generator of the replay files
(\S\ref{ssec:reproduce});
\item \texttt{papers/zeta/}: the sources and the PDF of this paper; \texttt{audit/followup/}: the logs and \texttt{\#print axioms} outputs of
the builds of the release.
\end{itemize}
The repository also contains the library \texttt{ZetaShell/} of a companion paper on families of Dirichlet $L$-functions, whose sources are in
\texttt{papers/family/}, and the library \texttt{ZetaQ/} added to the repository earlier (see \texttt{NOTICE}); \texttt{ZetaS} imports neither.
The statement files of~\cite{AF} (\texttt{Zeta23/Statement.lean} and the existing files of \texttt{comparator/}) are unchanged. Further,
the repository contains changes to $25$ files of the library \texttt{Zeta23} ($3{,}122$ lines added, $121$ deleted, between the upstream base
commit and \texttt{bb7223b}): a weaker parameter condition
(\texttt{Params.ValidQ} in place of \texttt{Params.Valid}), with the proofs adapted, and new modules for Gevrey ramps and a free buffer
(\texttt{Tail/GevreyTail}, \texttt{Taper/GevreyProduct}, \texttt{WindowD}). The axiom census above is for the tree with these changes.
The Lean files of \texttt{Zeta23}, \texttt{ZetaS}, \texttt{ZetaSCertShims} and \texttt{comparator/} in the release are those of commit
\texttt{bb7223b}, at which the census of \S\ref{app:census} was taken and against which the replay ran, except for comments: for the release,
the copyright headers of the four new comparator files of commit \texttt{bb7223b} and stale doc-strings in
\texttt{comparator/ChallengeDeps/FollowUpZeta.lean} and \texttt{ZetaS/ChallengeZetaS.lean} were corrected. Two files were added to
\texttt{comparator/PrintAxioms/}, \texttt{FollowUpReplay.lean} (\S\ref{app:census}) and \texttt{FollowUpShell.lean}, for the companion paper,
and \texttt{comparator/README\_followup.md} was brought up to date. The comparator package \texttt{FollowUpFamily} of the companion paper
(\texttt{ChallengeDeps/}, \texttt{Challenge/}, \texttt{Solution/} and \texttt{PrintAxioms/FollowUpFamily.lean}, \texttt{config-followup-family.json}
and \texttt{README\_followup\_family.md}) was also added.
The folders of \texttt{ZetaS}:
\begin{itemize}
\item \texttt{ZetaS/Interfaces.lean}, \texttt{InterfacesV2.lean}: the interfaces between the zero side and the local arguments (Gram data,
frame families at bandwidth $\lambda<1$); \texttt{ZetaS/ChallengeZetaS.lean}: the library's copies of the definitions of \S\ref{app:defs},
together with \texttt{MajorantCert} and \texttt{CertMaj};
\item \texttt{LinAlg/}: Section~\ref{sec:linalg} (stability inequality and its two-parameter form, $\Phi_m$, pinching, monotonicity, pair
removal, localisation, inertia);
\item \texttt{Window/}: the two windows, their exact moments and constants, and the window-general pipeline at bandwidth $\lambda<1$;
\texttt{Bandwidth/}: the limit $\lambda\to1^-$;
\item \texttt{ZeroSide/}: the Poisson and Gram identities, positive definiteness, truncation defect, inertia of pair blocks, and the assembly
from the zeros of \texttt{riemannZeta};
\item \texttt{SigmaDist/}: Section~\ref{sec:simple}; \texttt{KSide/}: Sections~\ref{sec:first} and~\ref{sec:sc}; \texttt{Top/}: the numerics of
the constants, the robustness conditions, the headline theorems (\texttt{TopFinal.lean}, \texttt{TopSSC.lean}) and the forms with a general
majorant threshold (\texttt{TopSolutionGeneral.lean});
\item \texttt{Majorant/}: the majorant, its analytic part, its cell checker, the $13$ replay shards and the two majorant theorems;
\texttt{Cert/}: the corrected checker for the $K=5$ certificate (\texttt{CheckerBase.lean}, \texttt{CheckerLeaves.lean},
\texttt{CheckerCoreV3.lean}, with kernel-checked negative controls in \texttt{CheckerTestsV3.lean}), the specification
\texttt{CertSpecAM5.lean}, the soundness nodes and the chain \texttt{ChainV3.lean};
\item \texttt{Skeleton/}: re-export files for the proved nodes, and the one open node \texttt{A8g\_SigmaDistGeneral.lean}.
\end{itemize}
The new Lean files of \texttt{comparator/} carry the author's copyright notice under the Apache License~2.0, the licence of the repository
of~\cite{AF}; the files that come from~\cite{AF} keep their own notices (see \texttt{NOTICE}).

